Null results in AI alignment research present a critical interpretational challenge:
when an intervention fails to induce measurable behavioral change, this can reflect
either genuine model robustness or experimental insufficiency arising from weak
interventions, small samples, or insensitive metrics. Classical null hypothesis
significance testing cannot distinguish these cases, yet alignment researchers
routinely interpret inconclusive null results as evidence of safety.

We develop and validate \ours, a Bayesian evidence quantification framework that
explicitly models intervention strength, measurement sensitivity, and prior effect
size distributions to yield a calibrated posterior over two competing hypotheses:
true robustness versus experimental weakness. The framework computes closed-form
Bayes factors (avoiding MCMC) and provides a sensitivity analysis over prior
hyperparameters to assess the robustness of conclusions.

Applied retrospectively to 27 published alignment null results spanning activation
steering, \rlhf, and jailbreak evaluations, \ours finds that \mainresult of null
results have posterior probability $P(\text{weakness} \mid \text{data}) > 0.3$---nearly
double our pre-registered threshold---and \emph{zero} published null results constitute
strong evidence of genuine robustness ($\mathrm{BF}_{01} > 10$). Prospective validation
using real \llama experiments reveals an unexpected ``inverse jailbreak'' effect:
weak educational-framing interventions achieve 20\% jailbreak success rates while
strong explicit jailbreak patterns yield 0\%, yet \ours correctly classifies the
strong-intervention null result as \emph{inconclusive} rather than evidence of
robustness.

These results demonstrate that alignment research is systematically producing
underpowered experiments whose null results mislead robustness assessment. We
release \ours as a practical diagnostic tool for the alignment community.
